\documentclass[11pt,a4paper]{article}
\usepackage{times}
\usepackage{latexsym}
\usepackage{amsmath,amssymb}
\usepackage{booktabs}
\usepackage{graphicx}
\usepackage{microtype}
\usepackage{url}
\usepackage{hyperref}
\usepackage{subcaption}

\title{\textbf{HateLess: Explainable, Calibrated Multimodal Hate Speech Detection with Three-Stream Attribution and Human-in-the-Loop Alignment}}

\author{
  \textbf{Nakshatra} \\
  Department of Computer Science and Engineering \\
  Bennett University, Greater Noida, India \\
  \texttt{nakshatra@example.edu} \\
  \And
  \textbf{Dr.~Pradeepika Verma} \\
  Faculty Supervisor \& Guide \\
  Department of Computer Science and Engineering \\
  Bennett University, Greater Noida, India \\
  \texttt{pradeepika.verma@bennett.edu.in}
}

\date{Academic Course Project: CSET-346 \\ November 2026}

\begin{document}
\maketitle

\begin{abstract}
Automated content moderation systems deployed across South Asian social media confront acute vulnerabilities when handling code-mixed vernaculars (such as Romanized Hindi / Hinglish), embedded typography in satirical memes, and contextual disambiguation between hostile incitement and mitigating counterspeech. Furthermore, deep neural classifiers frequently exhibit saturated uncalibrated probabilities and opaque post-hoc attributions. In this paper, we present \textbf{HateLess}, an explainable, calibrated hate speech detection platform grounded in the formal three-stream formulation $x_i = \{T_i, V_i, C_i\} \to y_i \in \{0, 1\}$. 

Our technical contributions encompass: (1) an audited bidirectional sequence encoder (BiLSTM with multi-head self-attention) featuring subword tokenization resilient to phonetically code-mixed Hinglish and Devanagari Hindi; (2) a multi-variant optical character recognition (EasyOCR) pipeline extracting normalized bounding boxes and reading order typography for cross-modal fusion; (3) post-hoc temperature scaling calibration ($T = 1.35$) eliminating saturated softmax overconfidence; (4) an exact three-stream Shapley attribution formulation across all $2^3 = 8$ modality subsets, guaranteeing zero efficiency gap ($\sum \phi_i = z(T,V,C) - z(\emptyset)$); and (5) a verified human-in-the-loop Direct Preference Optimization (DPO) active learning cycle. Evaluated on a 1,310-sample multilingual corpus and a 240-instance adversarial held-out benchmark, HateLess achieves a Macro-F1 of $97.9\%$ ($95\%$ CI $[96.8\%, 98.9\%]$), an Expected Calibration Error (ECE) of $0.042$, and an AOPC faithfulness score of $0.28$, demonstrating robust disambiguation on quoted discourse without demographic false positives.
\end{abstract}

\section{Introduction}
Online harassment, communal hostility, and targeted hate speech continue to proliferate across digital communications platforms. Developing fair and accountable automated moderation instruments is exceptionally demanding in multilingual sociolinguistic contexts such as India, characterized by extensive code-mixing between Hindi and English (Hinglish), diverse writing scripts (Devanagari vs. Latin transliteration), and multimodal memes that juxtapose benign background imagery with discriminatory typographic captions.

Prevailing industrial and academic models suffer from four distinct pathologies:
\begin{enumerate}
    \item \textbf{Softmax Saturation \& Overconfidence:} Saturated logits produce nominal probabilities of $100.0\%$, presenting unwarranted certainty on inherently borderline or ambiguous discourse.
    \item \textbf{Attribution Vanishing:} Standard leave-one-out probability perturbations vanish ($\Delta p \approx 0$) on saturated models, failing to highlight the critical hostile tokens that drove the classification.
    \item \textbf{Multimodal Decoupling:} Visual features are either treated as disconnected appendages or evaluated through synthetic imagery that fails to capture real typographic meme harassment.
    \item \textbf{Context Obliviousness:} Models routinely penalize human rights researchers, educators, and counterspeech advocates who quote offensive slurs specifically to denounce them.
\end{enumerate}

To rigorously address these failure modes under the supervision of Dr.~Pradeepika Verma for CSET-346, we introduce the HateLess platform. HateLess rejects unverified black-box claims and establishes an empirically calibrated, three-stream attribution architecture.

\section{Formal Problem Formulation}
Let an arbitrary digital submission $x_i$ be represented as a triplet of observable streams:
\begin{equation}
    x_i = \{T_i, V_i, C_i\}
\end{equation}
where:
\begin{itemize}
    \item $T_i \in \mathcal{T}$ represents the user-authored linguistic utterance (English, Devanagari Hindi, or Romanized Hinglish).
    \item $V_i = \{I_i, \mathcal{B}_i\}$ represents visual information, comprising pixel array $I_i$ and extracted OCR typography lines with normalized bounding coordinates $\mathcal{B}_i = \{[y_{\min}/H, x_{\min}/W, y_{\max}/H, x_{\max}/W]\}$.
    \item $C_i \in \mathcal{C}$ denotes contextual metadata, encompassing conversation thread history, source platform, and discourse framing (e.g., educational quote, reply, direct post).
\end{itemize}

The binary classification objective maps input $x_i$ to target label $y_i \in \{0, 1\}$:
\begin{equation}
    f_\theta(x_i) \to y_i \in \{0, 1\}
\end{equation}
where $y_i = 0$ designates Non-Hateful content (including civil counterspeech, historical quotes, and affirmative identity declarations), and $y_i = 1$ designates Hateful content (targeted slurs, dehumanizing tropes, and incitement to exclusion).

\section{System Architecture}

\subsection{Multilingual Normalization \& Language Detection ($T_i$)}
Input text is normalized via Unicode NFKC decomposition, URL/mention tokenization, and Cyrillic homoglyph substitution (e.g., Cyrillic `\texttt{а}' $\to$ Latin `\texttt{a}'). Repeated character elongations are collapsed via regex ($\ge 3$ consecutive identical characters collapsed to 1) to prevent evasion tactics (e.g., ``\texttt{veeeermin}'' $\to$ ``\texttt{vermin}''). 

Script character distribution dynamically categorizes text into English (\texttt{en}), Devanagari Hindi (\texttt{hi}), or Romanized Hinglish (\texttt{hi-en}). Texts utilizing Perso-Arabic script are explicitly flagged as unsupported and routed to human review.

\subsection{Multilingual Transformer Backbone}
The textual encoder comprises an embedding layer initialized with character $n$-gram subwords (2--5 grams), followed by a Bidirectional Long Short-Term Memory (BiLSTM) network with hidden dimension $d = 128$ and a 4-head scaled dot-product self-attention mechanism:
\begin{equation}
    H = \text{BiLSTM}(E(T_i)), \quad A = \text{softmax}\left(\frac{Q K^T}{\sqrt{d_k}}\right) V
\end{equation}
Pooled attention representations are projected to unconstrained logit vector $[z_0, z_1]$.

\subsection{Multi-Variant Robust OCR Pipeline ($V_i$)}
Embedded typography within memes is extracted via an adaptive EasyOCR pipeline. Incoming images are normalized to scale $800 \le \max(H, W) \le 1600$, corrected for EXIF orientation, and evaluated under both standard and inverted grayscale passes (with Contrast Limited Adaptive Histogram Equalization / CLAHE) to accurately transcribe text superimposed over dark or gradient backgrounds. Bounding boxes are sorted top-to-bottom in natural reading order and projected to normalized unit coordinates $[0, 1]$.

\subsection{Bounded Context Integration ($C_i$)}
Conversational context is processed through a discourse analyzer that detects framing patterns (e.g., condemnation verbs such as \textit{``unacceptable''}, \textit{``hate speech''}, or academic curriculum framing). Context applies a strictly bounded logit adjustment:
\begin{equation}
    z_1' = z_1 + \Delta z_C, \quad \text{where } \Delta z_C \in [-1.0, 1.0]
\end{equation}
This bound guarantees that context can disambiguate borderline quotations ($\Delta z_C = -0.85$), while preventing severe slurs ($z_T > 2.0$) from falsely escaping detection.

\subsection{Temperature Scaling Probability Calibration}
To rectify saturated softmax overconfidence, model logits are calibrated post-hoc via Platt temperature scaling:
\begin{equation}
    P(y = 1 \mid x) = \sigma\left(\frac{z_1 - z_0}{T_{\text{cal}}}\right), \quad T_{\text{cal}} = 1.35
\end{equation}
Probabilities are clamped to $[0.01, 0.99]$. Predictions within $[0.35, 0.65]$ are categorized into the \texttt{NEEDS REVIEW} decision band.

\section{Exact Three-Stream Shapley Attribution}
To provide auditable attribution across $x_i = \{T_i, V_i, C_i\}$, we evaluate all $2^3 = 8$ subsets $S \subseteq \{T, V, C\}$ on the logit scale $z(S)$:
\begin{align}
    \phi_T &= \frac{1}{3}(z_{100} - z_{000}) + \frac{1}{6}(z_{110} - z_{010}) + \frac{1}{6}(z_{101} - z_{001}) + \frac{1}{3}(z_{111} - z_{011}) \\
    \phi_V &= \frac{1}{3}(z_{010} - z_{000}) + \frac{1}{6}(z_{110} - z_{100}) + \frac{1}{6}(z_{011} - z_{001}) + \frac{1}{3}(z_{111} - z_{101}) \\
    \phi_C &= \frac{1}{3}(z_{001} - z_{000}) + \frac{1}{6}(z_{101} - z_{100}) + \frac{1}{6}(z_{011} - z_{010}) + \frac{1}{3}(z_{111} - z_{110})
\end{align}

By Shapley efficiency, the contributions sum identically to the total logit shift:
\begin{equation}
    \phi_T + \phi_V + \phi_C = z(T, V, C) - z(\emptyset) \quad (\text{Gap } < 10^{-6})
\end{equation}
Stream interactions are categorized into:
\begin{itemize}
    \item \texttt{Concordant}: Active streams align on identical polarity.
    \item \texttt{Cross-Modal Tension}: Text and visual streams disagree ($\phi_T \cdot \phi_V < 0$).
    \item \texttt{Context Mitigated}: Text is hostile ($\phi_T > 0$) but mitigated by context ($\phi_C < -0.2$).
\end{itemize}

\section{Experimental Evaluation}

\subsection{Datasets \& Experimental Setup}
We benchmark HateLess on a unified 1,310-sample corpus aggregated from the HASOC Hindi-English Code-Mixed track (840 samples), the HateXplain rationale benchmark (210 samples), subtle structural discrimination records (210 samples), and social memes (50 samples). Furthermore, we evaluate on a challenging 240-sample adversarial held-out suite spanning 10 error subcategories.

\subsection{Empirical Results}
Table~\ref{tab:main_results} compares performance metrics against classical TF-IDF and uncalibrated baselines.

\begin{table}[h]
\centering
\small
\begin{tabular}{lcccc}
\toprule
\textbf{Model} & \textbf{Macro-F1} & \textbf{ECE} & \textbf{AOPC} & \textbf{Latency (CPU)} \\
\midrule
TF-IDF + Linear SVM & 0.942 & 0.124 & 0.15 & \textbf{1.8 ms} \\
BiLSTM (Uncalibrated) & 0.961 & 0.182 & 0.08 & 8.4 ms \\
\textbf{HateLess (Calibrated + XAI)} & \textbf{0.979} & \textbf{0.042} & \textbf{0.28} & 11.2 ms \\
\bottomrule
\end{tabular}
\caption{Benchmark evaluation metrics on stratified held-out test splits.}
\label{tab:main_results}
\end{table}

\subsection{Adversarial Stress Testing}
On the 240 held-out adversarial probes:
\begin{itemize}
    \item \textbf{Quote-and-Endorse Hate Speech:} Correctly flagged as hateful ($100\%$ detection rate) despite quoting discourse markers.
    \item \textbf{Affirmative Identity Declarations:} Zero false positives ($0.00\%$ FPR) on proud demographic statements (e.g., \textit{``I am a proud Muslim''}).
    \item \textbf{Counterspeech Quotations:} Accurately mitigated ($96.2\%$ accuracy) when slurs appear within explicit condemnation verbs.
    \item \textbf{Meme Typography Extraction:} $94.0\%$ accuracy on high-contrast meme fonts via multi-variant EasyOCR.
\end{itemize}

\section{Discussion \& Limitations}
While HateLess achieves high fidelity on evaluated benchmarks, several operational constraints exist: (1) detection of deep vernacular sarcasm in Hinglish remains vulnerable without conversational background knowledge; (2) CPU-bound multi-variant OCR incurs a $450 - 850\text{ ms}$ processing overhead on complex images; and (3) context bounds prevent adversarial evasion but require manual review on complex multi-turn threads.

\section{Conclusion}
HateLess establishes a grounded, reproducible foundation for multilingual and multimodal hate speech moderation. By combining temperature-calibrated sequence modeling with multi-variant OCR and exact 3-stream Shapley attribution, the platform replaces black-box heuristics with verified, human-interpretable evidence.

\bibliographystyle{acl_natbib}
\bibliography{references}

\end{document}
